	\newpage
	\lstdefinestyle{bashstyle}{
		language=bash,
		basicstyle=\ttfamily\footnotesize,
		keywordstyle=\color{blue!90!black}\bfseries,
		commentstyle=\color{gray}\itshape,
		stringstyle=\color{green!70!black},
		numbers=left,
		numberstyle=\ttfamily\tiny\color{gray}\raisebox{0.5ex},
		stepnumber=1,
		numbersep=5pt,
		breaklines=true,
		frame=single,
		backgroundcolor=\color{gray!5},
		showspaces=false,
		showstringspaces=false,
		showtabs=false,
		tabsize=4,
		captionpos=b,
		linewidth=\linewidth,
		xleftmargin=20pt,
		xrightmargin=5pt,
		framexleftmargin=15pt,
		framexrightmargin=5pt,
		rulecolor=\color{blue!30},
		literate=  % Corrected Polish character mappings
		  {ą}{{\k{a}}}1 
		  {ć}{{\'c}}1 
		  {ę}{{\k{e}}}1 
		  {ł}{{\l}}1 
		  {ń}{{\'n}}1
		  {ó}{{\'o}}1 
		  {ś}{{\'s}}1 
		  {ź}{{\'z}}1 
		  {ż}{{\.{z}}}1
		  {Ą}{{\k{A}}}1 
		  {Ć}{{\'C}}1 
		  {Ę}{{\k{E}}}1 
		  {Ł}{{\L}}1
		  {Ń}{{\'N}}1 
		  {Ó}{{\'O}}1 
		  {Ś}{{\'S}}1 
		  {Ź}{{\'Z}}1 
		  {Ż}{{\.{Z}}}1,
		emph={sudo,apt-get,grep,awk,sed,chmod,chown,mv,cp,rm,mkdir,rmdir,cat,ps,kill,service,systemctl,apt,domain,ypinit,ypserv,ypbind,ypwhich,ypcat,ypmatch,ypdomainname},
		emphstyle=\color{magenta!70!black},
		lineskip=2pt,
		aboveskip=10pt,
		belowskip=10pt,
		numberblanklines=false
	}
\lstset{style=bashstyle} % Ustawienia dla całego dokumentu
\section{Procedury instalacyjne poszczególnych usług}		%4
% Procedury instalacyjne poszczególnych usług.
% (W podpunktach zamieścić polecenia dotyczące instalacji wdrażanych usług) 
\subsection{Setup Maszyn Wirtualnych}

Wszystkie poniższe kroki tj: \OznaczZdjecie[Rys.]{setup-01} - \OznaczZdjecie[Rys.]{setup-07} są wykonywane na maszynie każdej maszynie wirtualnej to jest: \texttt{MAIN}, \texttt{node1} oraz \texttt{node2}.

\fg[0.75\linewidth]{rys/setup/01.png}{Rozpoczęcie instalacji serwera}{setup-01}
\fg[0.75\linewidth]{rys/setup/02.png}{Wybór języka systemu}{setup-02}
\clearpage
\fg[0.75\linewidth]{rys/setup/03.png}{Wybór podstawy instalacji}{setup-03}
\fg[0.75\linewidth]{rys/setup/04.png}{Wybór adresu dla karty sieciowej}{setup-04}
\fg[0.75\linewidth]{rys/setup/05.png}{Wybór adresu mirroringu}{setup-05}
\fg[0.75\linewidth]{rys/setup/06.png}{Analiza przestrzeni dyskowej}{setup-06}
\fg[0.75\linewidth]{rys/setup/07.png}{Poprawna instalacja serwera}{setup-07}
\clearpage
Tutaj 2 opcjonalne kroki, które można wykonać w celu zainstalowania GUI na serwerze oraz dodatków gościa VirtualBox. Nie są one wymagane do działania usług NIS/NFS, ale mogą być przydatne w przypadku, gdy chcemy mieć dostęp do GUI na serwerze.

	\begin{lstlisting}[ caption={Instalacja GUI na serwerze}, label={lst:setup-gui}]
	
	sudo apt update && sudo apt upgrade -y

#(dla komfortu pracy.. kompletnie zbędny dla funkcjonalnosci)
	sudo apt install ubuntu-desktop 

	sudo systemctl set-default graphical.target

#zamontowac wczesniej addon..

	sudo mount /dev/cdrom /mnt
	cd /mnt
	sudo ./VBoxLinuxAdditions.run

#potem w opcjach (2 kierunkowy schowek w celu ułatwienia pracy z GUI)
	\end{lstlisting}
% \fg[0.75\linewidth]{rys/setup/08.png}{08}{setup-08}
% \fg[0.75\linewidth]{rys/setup/bruh2.png}{bruh2}{setup-bruh2}
% \fg[0.75\linewidth]{rys/setup/bruh.png}{bruh}{setup-bruh}

\subsection{Instalacja i konfiguracja NIS (main)}
\label{sec:nis-installation}
Rozpoczynamy od instalacji i konfiguracji NIS na serwerze głównym \texttt{MAIN}.\\
Najpierw aktualizujemy system i instalujemy pakiet NIS tak jak na \OznaczZdjecie[Rys.]{NIS-01}.
\fg[0.75\linewidth]{rys/NIS/01.png}{01}{NIS-01}
\clearpage
Następnie konfigurujemy adresy IP maszyn wirtualnych, aby mogły się ze sobą komunikować. W tym wchodzimy w ustawienia sieciowe Ubuntu i wybieramy opcję \texttt{Network} z menu. Wybieramy kartę sieciową, która ma być używana do komunikacji między maszynami wirtualnymi i ustawiamy jej adres IP na statyczny. Dla serwera \texttt{MAIN} możemy ustawić adres IP na \texttt{192.168.10.10} \OznaczZdjecie[Rys.]{NIS-02} i \OznaczZdjecie[Rys.]{NIS-03}. Dla pozostałych maszyn wirtualnych \texttt{node1} i \texttt{node2} ustawiamy adresy IP na \texttt{192.168.10.20} i \texttt{192.168.10.30} odpowiednio. Ustawiamy również maskę podsieci na \texttt{255.255.255.0}, brama oraz reszta ustawień pozostaje domyślna.\\

\fg[0.75\linewidth]{rys/NIS/02.png}{Konfiguracja ip \texttt{MAIN} cz.1}{NIS-02}
\clearpage
\fg[0.75\linewidth]{rys/NIS/03.png}{Konfiguracja ip \texttt{MAIN} cz.2}{NIS-03}
% \fg[0.75\linewidth]{rys/NIS/04.png}{04}{NIS-04}
Nastepnie ustawiamy hostname dla serwera \texttt{MAIN} na \texttt{nisArek}, oraz "stały" hostname w pliku \texttt{/etc/defaultdomain} na \texttt{nisArek}. Zmieniamy też uprawnienia na \texttt{700}.\\

\fg[0.75\linewidth]{rys/NIS/05.png}{Ustawianie nazwy domeny}{NIS-05}
\clearpage
Nastepnie edytujemy plik \texttt{/etc/default/nis} i zmieniamy w nim nazwę domeny na \texttt{nisArek} oraz status serwera na \texttt{master}.\\

\fg[0.75\linewidth]{rys/NIS/06.png}{Konfiguracja pliku \texttt{/etc/default/nis}}{NIS-06}

\vspace{-1cm}
Nastepnie konfigurujemy plik \texttt{etc/ypuserv.securenetes} tak aby zawierał sieci którym chcemy udostępnić nasz serwer NIS. W naszym przypadku jest to \texttt{192.168.10.0/24} oraz localhosty (ipv4 i ipv6)\\
\fg[0.75\linewidth]{rys/NIS/07.png}{Konfiguracja pliku \texttt{etc/ypuserv.securenetes}}{NIS-07}
\clearpage

Potem dodajemy linikę do pliku \texttt{/etc/yp.conf} aby nasz serwer NIS był widoczny w sieci:\\
\begin{lstlisting}[ caption={Dodanie linii do pliku \texttt{/etc/yp.conf}}, label={lst:nis-ypconf}]
	domain nisArek server 192.168.10.10
	\end{lstlisting}
% \fg[0.75\linewidth]{rys/NIS/08.png}{08}{NIS-08}
Tutaj sprawdzamy na jakich portach działa nasz serwer NIS.\\
\fg[0.75\linewidth]{rys/NIS/09.png}{Porty NIS}{NIS-09}
% \fg[0.75\linewidth]{rys/NIS/10.png}{10}{NIS-10}

Następnie aktywujemy uruchomianie serwera NIS przy starcie systemu.
\fg[0.75\linewidth]{rys/NIS/11.png}{Uruchomianie przy starcie}{NIS-11}
\clearpage
Potem uruchomiamy \texttt{/usr/lib/yp/ypinit -m} aby zainicjować listę użytkowników i grup.\\
\fg[0.75\linewidth]{rys/NIS/12.png}{Inicjacja \texttt{/usr/lib/yp/ypinit -m}}{NIS-12}

\fg[0.75\linewidth]{rys/NIS/13.png}{Powodzenie \texttt{/usr/lib/yp/ypinit -m}}{NIS-13}
\fg[0.75\linewidth]{rys/NIS/14.png}{Lista hostów, usług i grup NIS}{NIS-14}
\clearpage
Kończymy konfigurację serwera NIS uruchamiając go poleceniem \texttt{sudo systemctl restart ypserv}.\\
Potem otwieramy porty na których komunikuje się NIS dla naszej sieci (\texttt{192.168.10.0/25})\\
\fg[0.75\linewidth]{rys/NIS/15.png}{Poprawne uruchomienie usługi}{NIS-15}
\fg[0.75\linewidth]{rys/NIS/16.png}{Otwarcie portów na których komunikuje sie NIS}{NIS-16}
\clearpage

\subsection{Instalacja i konfiguracja NIS (node1 i node2)}

Poniższe kroki są wykonywane na maszynach klienckich \texttt{node1} oraz \texttt{node2}. Jeżeli nie jest to wyspecyifikowane inaczej, kroki są takie same jak dla \texttt{node1} i \texttt{node2}.\\

\label{sec:nis-installation-node}
Zaczynamy od instalacji NIS oraz ustawieniu hostname:\\
\fg[0.75\linewidth]{rys/NIS-CLIENT/01.png}{Instalacja usługi}{NIS-CLIENT-01}
\fg[0.75\linewidth]{rys/NIS-CLIENT/02.png}{Zmiana hostname}{NIS-CLIENT-02}

Potem dodajemy linikę do pliku \texttt{/etc/yp.conf} aby nasz serwer NIS był widoczny dla Node1 i Node2:\\
\fg[0.75\linewidth]{rys/NIS-CLIENT/03.png}{03}{NIS-CLIENT-03}

\clearpage
Następnie restartujemy usługę \texttt{ypbind}.\\

\fg[0.75\linewidth]{rys/NIS-CLIENT/04.png}{Reset usługi ypbind}{NIS-CLIENT-04}

Wykonujemy komendę \texttt{sudo reboot}, potem sprawdzamy zawartość pliku passwd, aby upewnić się, że użytkownicy są widoczni w systemie (W tym momencie nie byli stworzeni użytkownicy dla nodeów).\\ 

\fg[0.75\linewidth]{rys/NIS-CLIENT/06-poReboot.png}{Weryfikacja usługi}{NIS-CLIENT-06-poReboot}
\clearpage

\subsection{Instalacja i konfiguracja NFS (main)}
\label{sec:nfs-installation}

Zaczynam od instalacji NFS na serwerze głównym \texttt{MAIN}.\\
\fg[0.75\linewidth]{rys/NFS/01.png}{Instalacja usługi NFS}{NFS-01}

Następnie tworzymy katalogi, które będą udostępniane przez NFS. W tym przypadku są to \texttt{/srv/public} oraz \texttt{/srv/restricted}.\\
\fg[0.75\linewidth]{rys/NFS/02.png}{Tworzenie katalogów}{NFS-02}

Następnie tworzymy użytkowników i grupy, które będą miały dostęp do katalogów.\\

\begin{lstlisting}[ caption={Tworzenie użytkowników i grup}, label={lst:nfs-users}]

	sudo adduser --shell /bin/bash --home /home/nisuser1 nisuser1
	sudo adduser --shell /bin/bash --home /home/nisuser2 nisuser2
	
	sudo groupadd nisusers
	
	sudo usermod -aG nisusers nisuser1  
	sudo usermod -aG nisusers nisuser2 
	
	sudo chmod 777 /srv/public        # Pełne uprawnienia dla wszystkich
	sudo chmod 770 /srv/restricted    # Ograniczone uprawnienia
	sudo chown nobody:nogroup /srv/public
	sudo chown root:nisusers /srv/restricted 
	\end{lstlisting}
\clearpage
% \fg[0.75\linewidth]{rys/NFS/03.png}{03}{NFS-03}
Następnie edytujemy plik \texttt{/etc/exports} aby dodać katalogi, które będą udostępniane przez NFS dla poszczególnych adresów IP. W tym przypadku są to \texttt{/srv/public} oraz \texttt{/srv/restricted}. \texttt{/srv/restricted} będzie dostępny dla obydwu użytkowników, ale tylko \texttt{node2} będzie mógł zapisywać zawartość do tego katalogu (atrybut \texttt{rw}).\\
\fg[0.75\linewidth]{rys/NFS/04.png}{Edycja pliku \texttt{/srv/restricted}}{NFS-04}

Potem wykonujemy polecenie \texttt{sudo exportfs -arv} aby zaktualizować plik \texttt{/etc/exports}. Potem uruchamiamy usługę NFS oraz sprawdzamy czy działa.\\

\fg[0.75\linewidth]{rys/NFS/05.png}{Polecenie \texttt{sudo exportfs -arv}}{NFS-05}
\clearpage
\fg[0.75\linewidth]{rys/NFS/06.png}{Sprawdzanie statusu \texttt{nfs-server} oraz folderów eksportowanych}{NFS-06}

\fg[0.75\linewidth]{rys/NFS/07.png}{Sprawdzanie zamontowanych katalogów}{NFS-07}
% \fg[0.75\linewidth]{rys/NFS/08.png}{08}{NFS-08}
\clearpage

% \fg[0.75\linewidth]{rys/NFS/09.png}{09}{NFS-09}